\bookchapter{Глава 67. Доводка и выпуск — тесты, отладка и публикация в App Store}

\chapterintro{Приложение готово — и что дальше? В этой главе три вещи на пути «от запускается» к «можно передавать людям»: как быстро и эффективно отлаживать и ловить баги, как писать тесты против регрессий и как отправить приложение в App Store на скачивание всему миру. Публикация — не просто красивый ритуал: это водораздел между «работой» и «продуктом».}

\booksection{Три кита отладки: print, точки останова, предпросмотр}

\begin{center}
\begin{tabularx}{\linewidth}{|X|X|X|}
\hline
\textbf{Инструмент} & \textbf{Для чего} & \textbf{Как пользоваться} \tabularnewline
\hline
print + консоль & Смотреть журнал работы & В коде print("значение=\textbackslash{}(x)"), вывод смотрите в консоли внизу Xcode \tabularnewline
Точки останова & Остановить программу на нужной строке и посмотреть значения переменных & Щёлкните рядом с номером строки — появится синяя метка; дойдя до неё, программа замрёт — идите по шагам \tabularnewline
Предпросмотр SwiftUI + отладчик & Смотреть интерфейс и иерархию представлений вживую & Холст \#Preview + инспектор View Hierarchy во время работы приложения \tabularnewline
\hline
\end{tabularx}
\end{center}

\begin{center}\small\color{gray}Три кита отладки\end{center}

\begin{tcbinfo}{Совет: как читать сообщение о падении?}
При падении Swift консоль печатает «Fatal error» с причиной, а Xcode останавливается на сбойной строке и подсвечивает её. Историю падений на реальном устройстве смотрите в Xcode Window Organizer Crashes (нужен платный аккаунт). Ключевые слова для поиска берите в формате «первая фраза причины ошибки + swiftui» — точный вопрос рождает точный ответ.

\end{tcbinfo}

\booksection{Тесты: страховка для кода}

Ручная проверка ненадёжна: починили одну функцию — тихо сломалась соседняя. Тесты — это и «таблетка от сожалений», и «бронежилет». В iOS фреймворк XCTest, двух разновидностей:

\begin{center}
\begin{tabularx}{\linewidth}{|X|X|X|X|}
\hline
\textbf{Уровень} & \textbf{Что проверяет} & \textbf{Target} & \textbf{Скорость} \tabularnewline
\hline
Модульные тесты & Чистая логика (вычисления, проверки) & MyFirstAppTests & Секунды \tabularnewline
UI-тесты & Поведение настоящего интерфейса & MyFirstAppUITests & Минуты \tabularnewline
Ручное тестирование & Впечатления живого человека & —— & Нужно всегда \tabularnewline
\hline
\end{tabularx}
\end{center}

\begin{center}\small\color{gray}Уровни тестирования iOS (в шаблонном проекте уже предустановлены два Test Target)\end{center}

\begin{lstlisting}[language=Swift]
import XCTest

// Выносим логику вроде «строка суммы -> число» в чистую функцию
func parseAmount(_ input: String) -> Double? {
    guard let v = Double(input.trimmingCharacters(in: .whitespaces)),
          v > 0 else { return nil }
    return v
}

final class AmountTests: XCTestCase {
    func testEmptyInputReturnsNil() {
        XCTAssertNil(parseAmount(""))
    }
    func testNegativeReturnsNil() {
        XCTAssertNil(parseAmount("-5"))
    }
    func testNormalValueParses() {
        XCTAssertEqual(parseAmount("12.5"), 12.5)
    }
}
\end{lstlisting}

Запуск тестов: CmdU. Когда всё зелёное, менять код можно смело. Проверка сумм и расчёты сводок в книге учёта должны быть под прикрытием модульных тестов.

\booksection{Подпись и упаковка: ставим на приложение «печать»}

Каждое приложение iOS перед установкой обязано быть подписано — доказательство, что «пакет от вас и не подменён по дороге». Иерархия такова:

\begin{itemize}
  \item Этап учёбы: бесплатный Apple ID + автоподпись Xcode (Personal Team), работает на реальном устройстве, продлевать каждые 7 дней;
  \item Этап распространения: платный аккаунт разработчика (99 долларов в год) выпускает сертификаты и профили — без них ни в магазин, ни на долгую установку.
\end{itemize}

Упаковка: меню Xcode Product Archive (архивирование), после сборки в Organizer — Distribute App выбрать App Store Connect загрузить. Архивировать надо с целью Generic iOS Device (не симулятор) — новички сплошь и рядом забывают этот шаг.

\textbf{Осторожно: потерян сертификат — обновления встали}\par

Если сертификат подписи и закрытый ключ пропали, обновления для старого приложения выпускать больше нельзя — пользователю останется лишь «совершенно новое» приложение. Храните сертификат в надёжной резервной копии (экспорт .p12 из Связки ключей) и не теряйте пароль. К счастью, Apple предлагает автоподпись и облачное управление сертификатами (галочка Automatically manage signing в Xcode) — это снимает большую часть ручных граблей.

\booksection{Полный путь публикации в App Store}

Зарегистрировать аккаунт разработчика

Подготовить материалы для магазина

Загрузить сборку

Заполнить сведения о приватности и для проверки

Проверка Apple (1-7 дней)

Публикация

\begin{center}\small\color{gray}Рисунок 5: шесть шагов публикации в App Store\end{center}

\begin{itemize}
  \item Регистрация: developer.apple.com, личный аккаунт — 99 долларов в год;
  \item Материалы: иконка (1024×1024, без скруглений — Apple сама отрежет), скриншоты (размеры iPhone), описание, ссылка на политику конфиденциальности;
  \item Анкета конфиденциальности: честно перечислите, какие данные собираете, — книга учёта без сети может ответить «не собираем никаких данных»;
  \item Проверка: первое представление нередко ждёт 1–7 дней; отклонили — правьте по каждому пункту «замечаний». Причины отказа Apple обычно формулирует очень конкретно — просто делайте, как сказано;
  \item TestFlight: до публикации раздайте тестовую версию друзьям на закрытое тестирование (до 10 000 человек), соберите круг отзывов и лишь потом подавайте официально — сразу видно профессионала.
\end{itemize}

\booksection{Материалы для магазина: первое впечатление и есть количество загрузок}

Страница в магазине — витрина, на которую пользователь смотрит, «заходя в лавку». Первое впечатление делают три материала:

\begin{itemize}
  \item Иконка: 1024×1024, не отделывайтесь скриншотом. Не умеете в дизайн — сгенерируйте лаконичный знак ИИ-инструментом; годится лишь та, что узнаётся с одного взгляда даже уменьшенной на рабочем столе;
  \item Скриншоты: минимум 2 набора (ходовые размеры iPhone), показывайте самые главные сценарии; к каждому — одна фраза человеческим языком («Запись за секунду» в десять раз лучше «Эффективного решения для учёта финансов»);
  \item Описание: первая фраза объясняет «что это + для кого»; не громоздите «самый профессиональный, самый мощный» — искренность действует сильнее преувеличений.
\end{itemize}

\begin{tcbpractice}{Практика: один раз «выпустите» книгу учёта}
\booksection{Отчёты о падениях и мини-чеклист производительности}

Перед публикацией на эти два дела стоит потратить по полчаса:

\booksubsection{Мониторинг падений}

То, что у вас в тестах ничего не падало, не значит, что не будет падать у пользователей. Apple бесплатно даёт отчёты о падениях в Xcode Organizer (платный аккаунт видит журналы падений с реальных устройств); способ попроще — подключить сторонний SDK (например, Firebase Crashlytics). Публиковаться без мониторинга падений — всё равно что водить машину с закрытыми глазами.

\booksubsection{Четыре пункта самопроверки производительности}

\begin{enumerate}
  \item Списки через LazyVStack / List: длинные списки не выстраивайте обычным VStack в лоб — ленивая загрузка это первый урок производительности;
  \item Не пихайте в интерфейс картинки-оригиналы: берите AsyncImage (удалённые изображения) и системный кеш — оригинал на 5 МБ способен повесить старый телефон;
  \item Не делайте тяжёлую работу в главном потоке: сеть и большие файлы — через async/await;
  \item Замерьте Instruments: встроенный в Xcode анализатор производительности (Time Profiler, проверка утечек памяти) — сперва меряем, потом ускоряем, никаких «на глазок».
\end{enumerate}

(1) Добавьте 3 модульных теста к логике вроде parseAmount (CmdU — всё зелёное); (2) через точку останова своими руками поймайте специально устроенное падение и почините его; (3) соберите пакет через Product Archive и пригласите друзей на закрытое тестирование через TestFlight; (4) если остались силы, заведите аккаунт разработчика и по-настоящему отправьте книгу учёта в магазин. Первые три шага реально успеть сегодня вечером.

\end{tcbpractice}

\begin{tcbsummary}{Краткие итоги главы}
\begin{itemize}
  \item Три кита отладки: журнал print, точки останова, предпросмотр SwiftUI и инспектор представлений
  \item XCTest: модульные тесты (CmdU) + UI-тесты + ручное тестирование, шаблоны предустановлены
  \item Подпись: бесплатный аккаунт — автоподпись (продление раз в 7 дней); распространять может лишь платный
  \item Упаковка: Generic iOS Device + Archive + Distribute; сертификат обязательно в резервную копию
  \item Шесть шагов публикации: регистрация материалы загрузка анкета конфиденциальности проверка выход в магазин; сперва закрытое тестирование в TestFlight
\end{itemize}

\end{tcbsummary}
